\documentclass[12pt]{article}

\newcommand{\duedate}{-----}
\newcommand{\assignment}{Lecture 16 Summary: Diffusion Models}

% Change the following to your name and UNI.
\newcommand{\name}{Aksel Kretsinger-Walters, adk2164}
\newcommand{\email}{}

% NOTE: Defining collaborators is optional; to not list collaborators, comment out the
% line below. Maximum of two collaborators per problem set.
%\newcommand{\collaborators}{Vig Vigerton (\texttt{UNI}), Alice Bob (\texttt{UNI})}
% No collaborators on PS 0

\makeatletter
\def\input@path{{../}{../../}{../../../}} % add as many parents as you need
\makeatother
\input{pset_template_NNDL.tex} %% DO NOT CHANGE THIS LINE

\begin{document}
\psetheader %% DO NOT CHANGE THIS LINE

\section{VQ-VAE: Discrete Latent Codes}
Vector-Quantized VAEs map images to discrete latent tokens.
An encoder converts a $256\times 256$ image into a $32\times 32$ grid of discrete codes;
a decoder reconstructs the image from these codes.
The training objective includes:
\[
		\log p(x \mid q(z_e(x)))
		+ \| \mathrm{sg}[z_e(x)] - e \|_2^2
		+ \beta \| z_e(x) - \mathrm{sg}[e] \|_2^2 .
\]
(See slide 3 for diagram.)

\section{Diffusion Models: Intuition}
Diffusion models gradually corrupt data with Gaussian noise.
A neural network learns to reverse this noising process by predicting the noise added at each step.
After training, sampling proceeds by starting from pure noise and iteratively denoising.
(Visual examples on slide 4.)

\section{Formal Diffusion Process}
A forward noising Markov chain adds noise over $T$ steps:
\[
		q(x_t \mid x_{t-1}) = \mathcal{N}(x_t; \sqrt{1-\beta_t} x_{t-1}, \beta_t I),
\]
leading to $x_T$ approaching an isotropic Gaussian.
The reverse process is parameterized by a neural network:
\[
		p_\theta(x_{t-1} \mid x_t) = \mathcal{N}(x_{t-1}; \mu_\theta(x_t, t), \sigma_t^2 I).
\]
(See slide 5.)

\section{Training Objective}
Training maximizes a variational lower bound on $\log p_\theta(x_0)$.
When forward variances $\beta_t$ are fixed, the posterior $q(x_{t-1} \mid x_t, x_0)$ has no learnable parameters.
Using the reparameterization:
\[
		x_t = \sqrt{\bar{\alpha}_t}x_0 + \sqrt{1-\bar{\alpha}_t}\,\epsilon,
\]
the network is trained to predict the noise $\epsilon$.
The simplified loss (Ho et al.\ 2021) is:
\[
		\mathbb{E}_{t, x_0, \epsilon}\,\| \epsilon - \epsilon_\theta(x_t, t)\|_2^2 .
\]
(Slide 6.)

\section{Sampling and Architecture}
Sampling applies the learned reverse transitions:
\[
		x_{t-1} = \frac{1}{\sqrt{\alpha_t}}\Big(x_t -
		\frac{1-\alpha_t}{\sqrt{1-\bar{\alpha}_t}}\,\epsilon_\theta(x_t, t)\Big) + \sigma_t z.
\]
Most implementations use a U-Net because input and output shapes match.
(Slides 7--9.)

\section{Summary of Key Properties}
\begin{itemize}
		\item Diffusion is parameterized as a Markov chain; transitions are Gaussian.
		\item Forward variances are fixed or learned; reverse means are learned.
		\item Training optimizes a variational upper bound on $-\log p(x)$.
		\item KL terms simplify due to Gaussian and Markov assumptions.
		\item Predicting noise directly yields the most stable training.
		\item U-Net-like architectures fit naturally because $x_t$ and $x_{t-1}$ share dimensionality.
\end{itemize}
(See slide 10.)

\section{Cold Diffusion}
Cold Diffusion replaces Gaussian noise with arbitrary image degradations (e.g., blur, pixelation, masks).
The reverse model learns to invert these transformations without stochastic noise.
(Slide 11.)

\section{Stable Diffusion}
Two extensions:
\begin{enumerate}
		\item Perform diffusion in a latent space (via an autoencoder or VQ-VAE), offering major speed benefits.
		\item Condition the reverse process on text embeddings, using cross-attention between text and image latents.
\end{enumerate}
(Slides 12--13.)

\end{document}
